%Packages
%\usepackage{...}

%%%%%%%%%%%%%%%%
% Define some Maximes
\addMaxime[author=Cioran, % Author
fr=N'a de convictions que celui qui n'a rien approfondi., % First version
en=Has convictions the one who deepened nothing., % Second version
choose=fr{,}en, % Display fr and en version
ratio=0.6] % Ratio of \textwidth use for the maxime.

\addMaxime[author=Charles Darwin's brother Erasmus, % Author
fr=Finalement{,} le raisonnement � priori est si satisfaisant pour moi que si\maximeNewline les faits ne correspondent pas{,} mon sentiment est : tant pis pour les faits.,  % First version
en=In fact the � priori reasoning is so entirely satisfactory to me that if\maximeNewline the facts won't fit in{,} why so much the worse for the facts is my feeling., % Second version
cite=Erasmus_Darwin, % Citation key (use cite{})
choose=fr, % Display only fr version
ratio=0.77] % Ratio of \textwidth use for the maxime.


%%%%%%%%%%%%%%%%
% Define some Glossaries
\addAcronymAndGlossary{wcrt} % key
{WCRT} % Short form, singular
{WCRTs} % Short form, plural
{\textbf{W}orst \textbf{C}ase \textbf{R}esponse \textbf{T}ime} % Full form, singular
{\textbf{W}orst \textbf{C}ase \textbf{R}esponse \textbf{T}imes} % Full form, plural
{The \gls{wcrt} of a task is the maximum duration between the activation of the task and the moment it finishes its execution} % Glossary entry (automatically add the entry in Acronyms, even if you never use it, be careful.

\addAcronymAndGlossary{rt} % key
{RT} % Short form, singular
{} % No short form in plural 
{\textbf{R}eal-\textbf{T}ime} % Full form, singular
{} % No full form in plural
{} % No glossary entry

%%%%%%%%%%%%%%%%
% Define the bibliographie file
\addbibresource{Bibliography/Bibliography}

%%%%%%%%%%%%%%%%
% Define the informations for the Title Page
\phdAuthorFirstName{Pierre}
\phdAuthorLastName{Courbin}
\phdUniversityTitle{\includegraphics[width=1\textwidth]{MSTIC_LOGO} \par \vfill}
\phdUniversity{l'Universit{\'e} Paris-Est}
\phdSpeciality{Informatique}

\phdTitleEnglish{Scheduling Hard Sequential or Parallel Pre-emptive Real-Time Tasks upon Identical Multiprocessor platforms}
\phdTitleFrench{Ordonnancement de T{\^a}ches Temps R{\'e}el Dures Pr{\'e}emptives S{\'e}quentielles ou Parall{\`e}les sur plateformes mutiprocesseur identique}
\phdDate{vendredi 13 D�cembre 2013}

\phdAbstractFrench{Mon r�sum� en fran�ais}

\phdAbstractEnglish{Mon r�sum� en anglais}

\phdKeywordsEnglish{real-time, scheduling, multiprocessor, parallel, fork-join}
\phdKeywordsFrench{temps r�el, ordonnancement, multiprocesseur, parall�lle, fork-join}

\phdAddAdvisor{Laurent}{George}{Universit� Paris-Est (LIGM), France}
\phdPresident{Pascal}{Richard}{ENSMA (LIAS), France}
\phdAddReviewer{Alan}{Burns}{University of York (ARTIST), U.K}
\phdAddReviewer{Pascal}{Richard}{ENSMA (LIAS), France}
\phdAddExaminator{Jo\"el}{Goossens}{Universit{\'e} Libre de Bruxelles (PARTS), Belgium}
\phdAddExaminator{Yves}{Sorel}{INRIA Rocquencourt (AOSTE), France}

\phdAddAffiliationFirst{ECE_LOGO_EN}{PhD prepared at \\ \textbf{ECE Paris -- LACSC} \\ Laboratoire d'Analyse et Contr{\^o}le des Syst{\`e}mes Complexe \\ 37, Quai de Grenelle \\ CS 71520 \\ 75725 PARIS CEDEX 15}
\phdAddAffiliationSecond{LISSI_LOGO}{PhD in collaboration with \\ \textbf{UPEC -- LiSSi (EA 3956)} \\ Laboratoire Images, Signaux et Syst{\`e}mes Intelligents \\ Domaine Ch{\'e}rioux \\ 122 rue Paul Armangot \\ 94400 Vitry sur Seine}

\phdDedication{\`A mes parents,\phdNewline
\`A mes grands parents, encore l� ou d�j� partis,\phdNewline
\`A mes nombreux fr�res et s\oe urs,\phdNewline
\`A toute ma famille,\phdNewline
\`A tous mes amis,\phdNewline
\`A toutes les personnes que j'ai eu le privil�ge de croiser.\phdNewline\phdNewline
Parce que chaque personne rencontr�e, m�me bri�vement,\phdNewline est une occasion exceptionnelle de se re-d�couvrir\phdNewline et de s'�merveiller.\phdNewline
\vfill\phdNewline
\vfill\phdNewline
\vfill\phdNewline
To my parents,\phdNewline
To my grandparents, still there or already gone,\phdNewline
To my many brothers and sisters,\phdNewline
To all my family,\phdNewline
To all my friends,\phdNewline
To all the people who have graced my life.\phdNewline\phdNewline
Because each encounter, even a brief one,\phdNewline is an occasion to re-imagine oneself\phdNewline and to marvel.}

\phdPDFInfos %% Add all theses informations in my PDF metadatas

%%%%%%%%%%%%%%%%
% Define the chapterRectangle informations
\makeatletter
\if@chapterRectangle
	\setlength{\chapterRectangleWidth}{0.9cm} %% Chapter rectangle width
	\setlength{\chapterRectangleSectionWidth}{0.3cm} %% Section rectangle width
	\setlength{\chapterRectangleSubSectionWidth}{0.1cm} %% SubSection rectangle width
	\renewcommand\chapterRectangleRatioPageHeight{0.50} %% Ratio of the page height used to display rectangles
	
	\addChapterRectangleColor[chapter=1, color=red, section=3]
	\addChapterRectangleColorSection[chapter=1, section=3, subsection=2]
	\addChapterRectangleColor[chapter=2, color=green, section=0] %% Put 0 to section if you have only one and you don't want to display the gray rectangle of section
	
\fi
\makeatother


%%%%%%%%%%%%%%%%
% Configure hyperref
% Configuration of HyperLink Color
\definecolor{linkcol}{rgb}{0,0,0.4} 
\definecolor{citecol}{rgb}{0.5,0,0} 

\hypersetup{
	linkcolor=linkcol, %couleur des liens hypertextes internes
	citecolor=citecol, %couleur des liens pour les citations
	urlcolor=linkcol, %couleur des liens pour les url
}

%%%%%%%%%%%%%%%%
% Configure package listings
\lstset{
	language=XML,
	% list your attributes here
	morekeywords={xmlns,version,algoName,name,autoPath,path,fileName,deadline,distribution,minU,maxU,nbMinTasks,number,nbProcessors,type,Scale}
}

%%%%%%%%%%%%%%%%
% Others definitions

%Modifier la profondeur de num�rotation des sections etc
\setcounter{secnumdepth}{4} 
\setcounter{tocdepth}{4}

%D�finitions
\newcommand{\equals}{\stackrel{\mathrm{def}}{=}}
\newcommand{\task}{\tau}
\newcommand{\proc}{\pi}
\newcommand{\thop}{\operatorname{th}}
\newcommand{\suchthat}{\slash}
\newcommand{\itemStyle}[1]{\emph{\textbf{#1}}}

%Emplacement des images
\graphicspath{{./}{Images/}}
